% Ч1.3 — Проценты: процент от числа, число по проценту, процентное отношение
% Старая методичка: стр. 91–92 (старое задание 10)

\begin{tema}{Ч1.3}
\opr{Процент} — сотая часть: $1\%=\frac{1}{100}$, \ $p\%=\frac{p}{100}$.
Сначала решаем, \emph{что принято за 100\,\%}, — от этого зависит всё остальное.

\begin{pfTablica}{lc}
Что найти & Как\\
$p\%$ от числа $x$ & $x\cdot\dfrac{p}{100}$\\
Число, если $p\%$ от него равны $a$ & $a\cdot\dfrac{100}{p}$\\
Сколько процентов $a$ составляет от $b$ & $\dfrac{a}{b}\cdot100\%$\\
На сколько процентов $A$ больше $B$ & $\dfrac{A-B}{B}\cdot100\%$\\
На сколько процентов $B$ меньше $A$ & $\dfrac{A-B}{A}\cdot100\%$\\
\end{pfTablica}

\opr{Деление в отношении.} Чтобы разделить $S$ в отношении $a:b$, делим на $a+b$ частей:
\[
  S\cdot\frac{a}{a+b} \quad\text{и}\quad S\cdot\frac{b}{a+b}.
\]

\begin{primer}
120 в отношении $2:3$: одна часть $120:5=24$, получаем 48 и 72.
\end{primer}

\begin{oshibka}
«$A$ больше $B$ на $p\%$» и «$B$ меньше $A$ на $p\%$» — разные утверждения:
50 больше 40 на 25\,\%, а 40 меньше 50 на 20\,\%. За 100\,\% берётся то,
\emph{с чем} сравнивают.
\end{oshibka}
\end{tema}
